\begin{abstract}

In dieser Bachelorarbeit wird ein Plasma (MIPS-\RN{1}) System um einem standardisierten System- und Peripheriebus erweitert und von einem FPGA des Herstellers Xilinx auf ein FPGA des Herstellers Altera/Intel portiert. Das Plasma (MIPS-\RN{1}) System wird weiter ausgebaut, sodass eine Kommunikation mit neuen Hardware-Komponenten und IP-Cores von Altera/Intel möglich wird. Eine ARM-CPU auf dem Altera/Intel-FPGA und ein Linux-Betriebssystem werden konfiguriert. Für das Plasma-System werden eine Firmware und eine Beispielanwendung entworfen. 

Mit den in dieser Bachelorarbeit entwickelten Linux-Anwendungen und einem Treiber wird die Korrektheit der Portierung und Funktionsfähigkeit der Schnittstellen zwischen den ARM- und Plasma (MIPS-\RN{1}) Systemen gezeigt. 


%In diesem Dokument werden die Portierung eines Plasma (MIPS-I) Systems auf ein \gls{FPGA} des Herstellers Altera/Intel und Erweiterung um zusätzliche Komponenten, \todo{die eine einfachere Nutzung weiterer Hardware-Komponenten erlauben}, beschrieben. Besonderes Augenmerk wird auf die verwendeten Bussysteme, die Anbindung an einen \gls{ARM}-Prozessor und die Treiberentwicklung für Linux gelegt. 
%
%In this document the porting of a Plasma (MIPS-I) System to Altera/Intel \glspl{FPGA} and extension with additional components to ease the usage of further hardware components is described. The emphasis lies on used bus-systems, the connection to an \gls{ARM}-Processors and driver development for Linux. 
\end{abstract}